% !TEX program = xelatex
\documentclass[aspectratio=169,professionalfonts]{ctexbeamer}

\usetheme{Madrid}
\usecolortheme{seahorse}
\setbeamertemplate{navigation symbols}{}
\setbeamertemplate{footline}[frame number]

\usepackage{amsmath,amssymb,bm,mathtools}
\usepackage{booktabs}
\usepackage{array}
\usepackage{tabularx}
\usepackage{tikz}
\usetikzlibrary{positioning,calc}

\DeclareMathOperator{\Tr}{Tr}
\newcommand{\Frob}[2]{\left\langle #1, #2 \right\rangle_{\mathrm F}}
\newcommand{\Dfull}{N_{\mathrm{det}}}
\newcommand{\Na}{N_{\alpha}}
\newcommand{\Nb}{N_{\beta}}
\newcommand{\Nstr}{N_{\mathrm{str}}}
\newcommand{\Kpair}{N_{\mathrm{pair}}}
\newcommand{\Cmat}{C}

\title[Unique-Spin-String for VBSCF]{VBSCF 中的 Unique-Spin-String 矩阵化算法}
\subtitle{从 determinant-pair 求和到 unique alpha / beta 空间收缩}
\author{初版汇报稿}
\institute{xmvb-cpp project}
\date{\today}

\begin{document}

\begin{frame}
  \titlepage
\end{frame}

\begin{frame}{问题背景：旧算法的主拓扑}
  \small
  对于结构基矩阵元，传统 full-determinant 展开写成
  \[
  M_{IJ}
  =
  \sum_{p=1}^{\Dfull}
  \sum_{q=1}^{\Dfull}
  c_{Ip}\, c_{Jq}\, K(D_p, D_q),
  \qquad
  D_p = \bigl(\alpha_{i(p)}, \beta_{j(p)}\bigr).
  \]

  \vspace{0.4em}
  \begin{block}{核心瓶颈}
    外层索引是 determinant pair，因此主拓扑是
    \[
    O(\Dfull^2).
    \]
    即使单个 pair 的核函数被优化，若仍完整遍历 determinant pairs，整体复杂度主干不会变。
  \end{block}

  \begin{itemize}
    \item same-spin 的很多工作只依赖于 \(\alpha\)-pair 或 \(\beta\)-pair，本质上被重复计算了很多次。
    \item 因此真正需要改变的不是某个 kernel 常数，而是求和索引和数据表示。
    \item 我们的目标是把 determinant-space 的重复，重组到 unique-spin-string 空间中消掉。
  \end{itemize}
\end{frame}

\begin{frame}{重复来自哪里：determinant 是 alpha / beta 的组合}
  \small
  考虑一个最简单的 toy example：
  \[
  D_1=(\alpha_1,\beta_1),\quad
  D_2=(\alpha_1,\beta_2),\quad
  D_3=(\alpha_2,\beta_1),\quad
  D_4=(\alpha_2,\beta_2).
  \]

  \begin{columns}[T]
    \begin{column}{0.53\textwidth}
      \begin{block}{同一 full determinant 空间}
        \[
        \begin{array}{c|cc}
             & \beta_1 & \beta_2 \\ \hline
          \alpha_1 & D_1 & D_2 \\
          \alpha_2 & D_3 & D_4
        \end{array}
        \]
      \end{block}
      \[
      \Dfull = 4,\qquad \Na = 2,\qquad \Nb = 2.
      \]
    \end{column}
    \begin{column}{0.43\textwidth}
      \begin{block}{旧路径中的重复}
        \begin{itemize}
          \item full pairs: \(4^2 = 16\)
          \item unique \(\alpha\)-pairs: \(2^2 = 4\)
          \item unique \(\beta\)-pairs: \(2^2 = 4\)
        \end{itemize}
      \end{block}
      同一个 \((\alpha_i,\alpha_j)\) 或 \((\beta_k,\beta_l)\) kernel 会在大量 determinant pairs 中反复出现。
    \end{column}
  \end{columns}
\end{frame}

\begin{frame}{核心思想：引入 unique-spin-string 重标记}
  \small
  定义 determinant 到 unique spin-string 的映射
  \[
  u_{\alpha}(p)\in \{1,\dots,\Na\},
  \qquad
  u_{\beta}(p)\in \{1,\dots,\Nb\}.
  \]

  对于结构 \(I\)，不再把它看成 determinant coefficient vector，
  而是定义 unique alpha / beta 空间上的系数矩阵
  \[
  \Cmat_I(a,b)
  =
  \sum_{p:\,u_{\alpha}(p)=a,\;u_{\beta}(p)=b}
  c_{Ip}.
  \]

  \begin{block}{算法对象发生了本质变化}
    旧对象：\(\{c_{Ip}\}_{p=1}^{\Dfull}\) \\
    新对象：\(\Cmat_I \in \mathbb{R}^{\Na \times \Nb}\)
  \end{block}

  \begin{itemize}
    \item determinant-pair 求和被改写为 unique-spin 空间中的矩阵乘法与 Frobenius 收缩。
    \item 这不是“给旧算法加 cache”，而是把结构基展开重新组织到了新坐标系。
  \end{itemize}
\end{frame}

\begin{frame}{same-spin 矩阵化公式：forward}
  \small
  当某个通道可写成 \(\alpha\) 与 \(\beta\) 的可分离核时，
  \[
  K(D_p,D_q)
  =
  A_{u_{\alpha}(p),u_{\alpha}(q)}
  B_{u_{\beta}(p),u_{\beta}(q)},
  \]
  则结构矩阵元可重写为
  \[
  M_{IJ}
  =
  \Frob{\Cmat_I}{A\, \Cmat_J\, B^\mathsf{T}}.
  \]

  \vspace{0.4em}
  对应到 VBSCF 的几个主要通道：
  \[
  S_{IJ}
  =
  \Frob{\Cmat_I}{S^\alpha \Cmat_J (S^\beta)^\mathsf{T}},
  \]
  \[
  H^{(1)}_{IJ}
  =
  \Frob{\Cmat_I}{
    H^{\alpha,(1)} \Cmat_J (S^\beta)^\mathsf{T}
    +
    S^\alpha \Cmat_J (H^{\beta,(1)})^\mathsf{T}},
  \]
  \[
  H^{\mathrm{ss}}_{IJ}
  =
  \Frob{\Cmat_I}{
    H^{\alpha,\mathrm{ss}} \Cmat_J (S^\beta)^\mathsf{T}
    +
    S^\alpha \Cmat_J (H^{\beta,\mathrm{ss}})^\mathsf{T}}.
  \]

  \begin{block}{结论}
    determinant-pair 外层双重求和被矩阵收缩替代，同一个 unique same-spin pair 只计算一次。
  \end{block}
\end{frame}

\begin{frame}{backward 也闭合在 unique-spin 空间}
  \small
  对于选中的状态 \(n\)，先构造其 unique-spin 系数矩阵
  \[
  \Cmat^{(n)} \in \mathbb{R}^{\Na\times\Nb}.
  \]

  same-spin backward 的权重矩阵可直接写成
  \[
  W^\alpha
  =
  \sum_n w_n \Cmat^{(n)} B \bigl(\Cmat^{(n)}\bigr)^\mathsf{T},
  \qquad
  W^\beta
  =
  \sum_n w_n \bigl(\Cmat^{(n)}\bigr)^\mathsf{T} A \Cmat^{(n)}.
  \]

  \begin{itemize}
    \item 不再从 generalized eigensolve 后回到 full determinant-pair adjoint。
    \item 先在 unique-spin 空间聚合 state-average 权重，再做 same-spin exact backward。
    \item 这也是为什么结构矩阵构建和 active-space gradient 都能一起加速。
  \end{itemize}

  \begin{block}{意义}
    我们完成的不是单纯 forward 改写，而是 forward / backward 一体化的矩阵化。
  \end{block}
\end{frame}

\begin{frame}{这项工作的本质是什么}
  \small
  \begin{block}{不是}
    \begin{itemize}
      \item 不是缓存 full determinant pair。
      \item 不是额外增加一个 scalar cache 来避免少量重复。
      \item 不是只在代码层面把某些循环做小修小补。
    \end{itemize}
  \end{block}

  \begin{block}{而是}
    \begin{itemize}
      \item 提出 unique-spin-string / unique-spin-pair 表示；
      \item 将 structure basis 与 selected states 都重写为 \(\Na\times\Nb\) 矩阵对象；
      \item 把 same-spin 的 forward 与 backward 从 determinant-pair topology 中剥离出来；
      \item 将热点计算转化为 BLAS 友好的 dense/symmetric matrix contractions。
    \end{itemize}
  \end{block}

  \[
  \boxed{
    \text{这是算法重排，而不是 cache trick。}
  }
  \]
\end{frame}

\begin{frame}{标度分析：same-spin 的本质改进}
  \small
  旧算法中，可复用的 same-spin kernel 仍被嵌在 determinant-pair 外层：
  \[
  O(\Dfull^2).
  \]

  新算法中，需要真正求值的 same-spin kernel 数量变为
  \[
  \Na^2 + \Nb^2.
  \]

  若 full determinant 空间近似 Cartesian product，
  \[
  \Dfull \approx \Na \Nb,
  \]
  则 same-spin 可复用部分从
  \[
  O(\Na^2 \Nb^2)
  \]
  下降到
  \[
  O(\Na^2 + \Nb^2).
  \]

  \begin{block}{10698 / 400 determinant 的例子}
    \[
    \Dfull = 400,\qquad \Na = \Nb = 20.
    \]
    需要处理的 same-spin kernel 量级从约
    \[
    1.6\times 10^5
    \]
    降到
    \[
    800.
    \]
  \end{block}
\end{frame}

\begin{frame}{当前 opposite-spin 的位置}
  \small
  opposite-spin 已经不再回到 full determinant-pair sweep，但它也还没有达到
  \[
  O(\Na^2 + \Nb^2)
  \]
  的极致分离。

  \begin{block}{当前 production 路径的理解}
    \begin{itemize}
      \item pair topology 已经换成 unique \(\alpha\)-pairs / \(\beta\)-pairs；
      \item 但 opposite-spin 本质上仍是 \(\alpha\)-\(\beta\) mixed contraction；
      \item 当前主标度更接近
      \[
      O\!\left(K(\Na^2\Nb + \Na\Nb^2)\right),
      \]
      其中 \(K\) 为 active packed-pair 数。
    \end{itemize}
  \end{block}

  \begin{itemize}
    \item 因此，same-spin 的 \(N_{\det}^2\) 已经被真正消掉。
    \item 当前剩余 frontier 是 opposite-spin 的进一步矩阵化、低秩化或 rank-space 闭合。
  \end{itemize}
\end{frame}

\begin{frame}{exact 与 RI 的统一理解}
  \small
  \begin{itemize}
    \item unique-spin-string 并不是只适用于某一条积分路径。
    \item 它的核心是：先在 unique-spin 空间组织 \(\Cmat_I\) 与 \(\Cmat^{(n)}\)，
      再把可复用的 same-spin / projected opposite-spin payload 组织成矩阵对象。
  \end{itemize}

  \begin{block}{exact 路径}
    same-spin kernel 直接来自 exact determinant overlap / Hamiltonian evaluator，
    然后进入 matrix-form forward/backward。
  \end{block}

  \begin{block}{RI 路径}
    RI 让 active-space 2e kernel 更便宜，因此 same-spin determinant-pair
    这部分更容易成为主瓶颈。此时 unique-spin-string 的收益会更明显地体现在端到端单步时间上。
  \end{block}

  \begin{alertblock}{当前状态}
    我们已经把 same-spin 算法完整打通到 exact 与 RI production 路径；
    opposite-spin 仍有进一步提升空间，但这不影响 same-spin 算法本身已经构成实质性改进。
  \end{alertblock}
\end{frame}

\begin{frame}{正确性：不是近似，而是严格重组}
  \small
  \begin{block}{数学性质}
    unique-spin-string 路径是 exact algebraic reformulation，不改变任何目标矩阵元与梯度定义。
  \end{block}

  \begin{itemize}
    \item \texttt{check\_structure\_builder\_fast\_path}, \texttt{10698}, exact:
    \[
    \max |\Delta S| \approx 2.22\times 10^{-16},
    \quad
    \max |\Delta H| \approx 1.78\times 10^{-15}.
    \]
    \item active-space two-electron gradient finite difference:
    \[
    \text{误差量级 } 10^{-8}\sim 10^{-9}.
    \]
    \item orbital gradient finite difference:
    \[
    \text{误差量级 } 10^{-9}.
    \]
  \end{itemize}

  \begin{block}{结论}
    这说明我们做的是严格等价的矩阵化重排，而不是数值近似或启发式截断。
  \end{block}
\end{frame}

\begin{frame}{benchmark：structure build 的直接收益}
  \small
  以 \texttt{10698} 为例，expanded determinant 数为 400，unique alpha / beta 数均为 20。

  \vspace{0.5em}
  \begin{tabular}{lccc}
    \toprule
    指标 & 旧路径 / reference & fast path & 加速比 \\
    \midrule
    exact structure build & \(0.2373\) s & \(0.0302\) s & \(7.9\times\) \\
    \bottomrule
  \end{tabular}

  \vspace{0.8em}
  从已有 kernel benchmark 看，same-spin unique-spin-string 在 production 场景下的收益更明显：

  \vspace{0.4em}
  \begin{tabular}{lccc}
    \toprule
    路径 & structure build 关闭算法 & 开启算法 & 加速比 \\
    \midrule
    exact & \(0.236\) s & \(0.00962\) s & \(24.6\times\) \\
    RI & \(9.869\) s & \(0.0632\) s & \(156\times\) \\
    \bottomrule
  \end{tabular}

  \vspace{0.5em}
  \begin{block}{解释}
    当 active-space 预处理本身很重时，总单步时间不一定同步下降同样倍数；但 structure matrix 与 active-space adjoint 的热点已经被显著压缩。
  \end{block}
\end{frame}

\begin{frame}{benchmark：端到端单步时间}
  \small
  对 \texttt{10698} 的 RI production 单步 benchmark，steady-state 的端到端表现为：

  \vspace{0.4em}
  \begin{tabular}{lccc}
    \toprule
    指标 & 关算法 & 开算法 & 加速比 \\
    \midrule
    same-spin-phi-cache + structure & \(10.701\) s & \(0.0228\) s & \(469\times\) \\
    active adjoint & \(20.772\) s & \(0.1769\) s & \(117\times\) \\
    active space total & \(32.172\) s & \(0.8277\) s & \(38.9\times\) \\
    total single-step & \(32.829\) s & \(1.444\) s & \(22.7\times\) \\
    \bottomrule
  \end{tabular}

  \vspace{0.6em}
  \begin{block}{汇报时建议强调}
    \begin{itemize}
      \item pair-count 的理论降幅并不会 1:1 变成总墙钟加速；
      \item 但在 RI 路径下，same-spin 的确已经成为决定单步时间的关键环节；
      \item 一旦把这个环节矩阵化，收益会真实反映到 production 端到端时间。
    \end{itemize}
  \end{block}
\end{frame}

\begin{frame}{汇总结论}
  \small
  \begin{enumerate}
    \item 我们提出了 unique-spin-string / unique-spin-pair 表示，将 determinant expansion 重新组织到 unique alpha / beta 空间。
    \item 我们消除了 same-spin determinant-pair 主干中的重复，forward 与 backward 都已矩阵化。
    \item 该方法是 strict reformulation，不改变矩阵元与梯度，数值一致到机器精度。
    \item 在 exact 与 RI production 路径中，结构矩阵构建和 active-space gradient 都得到显著加速。
    \item 当前进一步优化的 frontier 已经从 same-spin 转移到 opposite-spin mixed contraction。
  \end{enumerate}

  \vspace{0.8em}
  \[
  \boxed{
  \text{Unique-spin-string 的价值，在于从算法层面把 same-spin 重复从 determinant 空间剥离出来。}
  }
  \]
\end{frame}

\begin{frame}{下一步可扩展的方向}
  \small
  \begin{itemize}
    \item 补充一页更细的 exact / RI opposite-spin 标度分析。
    \item 加入 one toy example 的 coefficient-matrix 动画或逐步示意图。
    \item 用实际 profile 图替换当前 benchmark 表。
    \item 如果用于正式组会，可增加 “代码实现映射” 页，标出
      \texttt{full\_structure\_builder.cpp}、\texttt{same\_spin\_matrix\_backward.cpp}、
      \texttt{opposite\_spin\_matrix\_backward.cpp} 三个落地点。
    \item 如果后续需要投稿导向版汇报，可直接扩展为
      ``motivation $\rightarrow$ theory $\rightarrow$ implementation
      $\rightarrow$ benchmark $\rightarrow$ outlook'' 五段式。
  \end{itemize}
\end{frame}

\end{document}
